%--- Part 6: Conclusions (3 min) -----------------------------------------------

\section{Conclusions}

%--- Slide 23 ------------------------------------------------------------------
\begin{frame}{Summary of Findings}
  \begin{enumerate}
    \item<1-> \textbf{RQ1 — Which characteristics drive prediction error?}
    \begin{itemize}
      \item \texttt{std\_speed} (speed variation) and
            \texttt{heading\_change\_per\_sec} (turning) are the
            dominant drivers in both GAM and XGBoost
      \item Scene-level features (\texttt{scene\_density\_VEHICLE}) are
            consistently least important; Trajectron++ handles scene
            context better than individual motion dynamics
    \end{itemize}
    \item<2-> \textbf{RQ2 — What are the success and failure modes?}
    \begin{itemize}
      \item Easy: slow or near-stationary pedestrians with low speed
            variation — Trajectron++ excels here
      \item Hard: continuous high-variation regime driven by erratic,
            fast-turning pedestrians; no sharp sub-clusters
    \end{itemize}
    \item<3-> \textbf{RQ3 — How do model settings affect performance?}
    \begin{itemize}
      \item Prediction horizon (\texttt{prediction\_sec}) has the
            largest effect — longer horizon raises error strongly
      \item Attention radius (\texttt{attention\_radius\_m}) is
            second; history window (\texttt{history\_sec}) has
            minimal impact
    \end{itemize}
  \end{enumerate}
\end{frame}

%--- Slide 24 ------------------------------------------------------------------
\begin{frame}{Limitations \& Future Work}
  \textbf{Limitations}
  \begin{itemize}[<+->]
    \item Analysis restricted to \emph{pedestrian} trajectories on
          nuScenes — generalisability to other agent types or datasets
          is not guaranteed
    \item Meta-models explain Trajectron++ predictions, not ground-truth
          difficulty — effects are model-specific
    \item Hard-regime clustering is weak (75\% noise); discrete
          failure-mode identification remains an open challenge
    \item Settings sweep uses nuScenes \emph{mini} —
          results may not transfer to full-scale training
  \end{itemize}
  \vspace{0.4em}
  \onslide<5->{%
  \textbf{Future Work}
  \begin{itemize}[<5->]
    \item Extend pipeline to vehicle, cyclist, and motorcyclist trajectories
    \item Apply to other trajectory prediction models and datasets
          via the trajdata interface
    \item Explore higher-capacity meta-models to better isolate
          failure modes in the hard regime
  \end{itemize}}
\end{frame}
